\chapter{Differential forms in Euclidean space. First order differential equations}

The equation of particular interest in this chapter is the first order equations of the form:
\begin{equation}
  \odv{y}{x} = f(x, y).
\end{equation}
In most texts

\section{Separable equations}

The R.H.S. of this equation might be tangled with \(y\), which makes it hard to solve. For some differential equations, the function in the R.H.S. can be written as a product of two functions: one dependent only on \(x\) and one depedendent only on \(y\):
\begin{equation}
  f(x, y) = X(x)Y(y)
\end{equation}
If this is the case, then the equation is \textbf{separable}, and solving it is easy.

\begin{df}{Separable differential equation}{}
  A first order differential equation is separable if it can be written in the form
  \begin{equation}
    \odv{y}{x} = X(x)Y(y)
  \end{equation}
  where \(X\) and \(Y\) are functions that are dependent only on \(x\) and \(y\) respectively.
\end{df}
